% GENERATED FILE. Edit canonical lessons, then run npm run generate:books.
% canonical-source-hash: fnv1a64:2296d151fb484328
% canonical-lessons: SA-C287-anudhavati, SA-C287-vegah, SA-C287-upasarpati, SA-C287-jagarti, SA-C287-kolahalah

\chapter{The Chase}
\label{ch:sa-the-chase}
\hlchaptermodality{287}

\begin{chapteropening}
\textbf{By the end of this chapter:} \emph{I can say who ran after the thief and what would have stopped him.}

Complete the last lesson of chapter 287: I can say who ran after the thief and what would have stopped him.
\end{chapteropening}

\section[\texorpdfstring{\sk{अनुधावति} (anudhāvati)}{अनुधावति (anudhāvati)}]{\sk{अनुधावति} (anudhāvati) — runs after; chases}

\label{lesson:SA-C287-anudhavati}



\begin{quote}
Before the new one: say the Sanskrit for a witness, then the Sanskrit for runs away.
\end{quote}

\subsection*{You'll want to know: \sk{अनुधावति}}
\textbf{\sk{अनुधावति}} — \emph{anudhāvati} — \textquotedblleft{}runs after; chases\textquotedblright{}.

\begin{quote}
\sk{यामिकः} \sk{चौरम्} \sk{अनुधावति।}
\end{quote}

\begin{quote}
\emph{yāmikaḥ cauram anudhāvati.}
\end{quote}

\begin{quote}
The watchman runs after the thief.
\end{quote}

Past (he …): \sk{सः} \sk{चौरम्} \sk{अनुधावितवान्}, \emph{saḥ cauram anudhāvitavān}.

\begin{cousinweb}
\textbf{\sk{अनुधावति}} (\emph{anudhāvati}) is built from \emph{anu}, \textquotedblleft{}after\textquotedblright{}, and the root √dhāv, \textquotedblleft{}to run\textquotedblright{}, which also gives \emph{dhāvati}, \textquotedblleft{}runs\textquotedblright{}; MW gives \emph{anu-dhāv} as to run after or pursue (MW).
\end{cousinweb}

\begin{grammarlens}[title={What you've built}]
One more word for telling what happened.
\end{grammarlens}

\subsection*{Your turn}
\begin{itemize}
\raggedright
  \item \emph{Say it:} \emph{anudhāvati}
  \item \emph{Say it:} \emph{anudhāvati}, once more
  \item \emph{Say it:} \emph{yāmikaḥ cauram anudhāvati.} — The watchman runs after the thief.
  \item \emph{Recall:} say the Sanskrit for a witness, then the Sanskrit for runs away, then say \emph{anudhāvati} again
\end{itemize}

\begin{tcolorbox}[breakable,colback=teal!4,colframe=teal!35!black,title={Before you move on}]
What does \sk{अनुधावति} mean? (\textquotedblleft{}Runs after; chases\textquotedblright{}.) Say it once more.
\end{tcolorbox}
\section[\texorpdfstring{\sk{वेगः} (vegaḥ)}{वेगः (vegaḥ)}]{\sk{वेगः} (vegaḥ) — speed}

\label{lesson:SA-C287-vegah}



\begin{quote}
Before the new one: say the Sanskrit for runs away, then the Sanskrit for runs after.
\end{quote}

\subsection*{You'll want to know: \sk{वेगः}}
\textbf{\sk{वेगः}} — \emph{vegaḥ} — \textquotedblleft{}speed\textquotedblright{}.

\begin{quote}
\sk{यामिकः} \sk{वेगेन} \sk{धावति।}
\end{quote}

\begin{quote}
\emph{yāmikaḥ vegena dhāvati.}
\end{quote}

\begin{quote}
The watchman runs at speed.
\end{quote}

A masculine noun.

\begin{cousinweb}
\textbf{\sk{वेगः}} (\emph{vegaḥ}) comes from the root √vij, \textquotedblleft{}to move with a quick, darting motion\textquotedblright{}; MW gives \emph{vega} as a violent rush and also as speed or velocity (MW).
\end{cousinweb}

\begin{grammarlens}[title={What you've built}]
One more word for telling what happened.
\end{grammarlens}

\subsection*{Your turn}
\begin{itemize}
\raggedright
  \item \emph{Say it:} \emph{vegaḥ}
  \item \emph{Say it:} \emph{vegaḥ}, once more
  \item \emph{Say it:} \emph{yāmikaḥ vegena dhāvati.} — The watchman runs at speed.
  \item \emph{Recall:} say the Sanskrit for runs away, then the Sanskrit for runs after, then say \emph{vegaḥ} again
\end{itemize}

\begin{tcolorbox}[breakable,colback=teal!4,colframe=teal!35!black,title={Before you move on}]
What does \sk{वेगः} mean? (\textquotedblleft{}Speed\textquotedblright{}.) Say it once more.
\end{tcolorbox}
\section[\texorpdfstring{\sk{उपसर्पति} (upasarpati)}{उपसर्पति (upasarpati)}]{\sk{उपसर्पति} (upasarpati) — creeps up to}

\label{lesson:SA-C287-upasarpati}



\begin{quote}
Before the new one: say the Sanskrit for runs after, then the Sanskrit for speed.
\end{quote}

\subsection*{You'll want to know: \sk{उपसर्पति}}
\textbf{\sk{उपसर्पति}} — \emph{upasarpati} — \textquotedblleft{}creeps up to\textquotedblright{}.

\begin{quote}
\sk{चौरः} \sk{शनैः} \sk{गृहम्} \sk{उपसर्पति।}
\end{quote}

\begin{quote}
\emph{cauraḥ śanaiḥ gṛham upasarpati.}
\end{quote}

\begin{quote}
The thief creeps slowly up to the house.
\end{quote}

Past (he …): \sk{सः} \sk{गृहम्} \sk{उपसृप्तवान्}, \emph{saḥ gṛham upasṛptavān}.

\begin{cousinweb}
\textbf{\sk{उपसर्पति}} (\emph{upasarpati}) is built from \emph{upa}, \textquotedblleft{}towards\textquotedblright{}, and the root √sṛp, \textquotedblleft{}to creep\textquotedblright{}; MW gives \emph{upa-sṛp} as to creep towards or approach softly (MW).
\end{cousinweb}

\begin{grammarlens}[title={What you've built}]
One more word for telling what happened.
\end{grammarlens}

\subsection*{Your turn}
\begin{itemize}
\raggedright
  \item \emph{Say it:} \emph{upasarpati}
  \item \emph{Say it:} \emph{upasarpati}, once more
  \item \emph{Say it:} \emph{cauraḥ śanaiḥ gṛham upasarpati.} — The thief creeps slowly up to the house.
  \item \emph{Recall:} say the Sanskrit for runs after, then the Sanskrit for speed, then say \emph{upasarpati} again
\end{itemize}

\begin{tcolorbox}[breakable,colback=teal!4,colframe=teal!35!black,title={Before you move on}]
What does \sk{उपसर्पति} mean? (\textquotedblleft{}Creeps up to\textquotedblright{}.) Say it once more.
\end{tcolorbox}
\section[\texorpdfstring{\sk{जागर्ति} (jāgarti)}{जागर्ति (jāgarti)}]{\sk{जागर्ति} (jāgarti) — stays awake; keeps watch}

\label{lesson:SA-C287-jagarti}



\begin{quote}
Before the new one: say the Sanskrit for speed, then the Sanskrit for creeps up to.
\end{quote}

\subsection*{You'll want to know: \sk{जागर्ति}}
\textbf{\sk{जागर्ति}} — \emph{jāgarti} — \textquotedblleft{}stays awake; keeps watch\textquotedblright{}.

\begin{quote}
\sk{यामिकः} \sk{सर्वां} \sk{रात्रिं} \sk{जागर्ति।}
\end{quote}

\begin{quote}
\emph{yāmikaḥ sarvāṃ rātriṃ jāgarti.}
\end{quote}

\begin{quote}
The watchman stays awake all night.
\end{quote}

Past (he …): \sk{सः} \sk{रात्रौ} \sk{जागरितवान्}, \emph{saḥ rātrau jāgaritavān}.

\begin{cousinweb}
\textbf{\sk{जागर्ति}} (\emph{jāgarti}) comes from the root √jāgṛ, \textquotedblleft{}to be awake, watch\textquotedblright{}; MW compares Greek \emph{egeirō} and Latin \emph{vigilo} (MW).
\end{cousinweb}

\begin{grammarlens}[title={What you've built}]
One more word for telling what happened.
\end{grammarlens}

\subsection*{Your turn}
\begin{itemize}
\raggedright
  \item \emph{Say it:} \emph{jāgarti}
  \item \emph{Say it:} \emph{jāgarti}, once more
  \item \emph{Say it:} \emph{yāmikaḥ sarvāṃ rātriṃ jāgarti.} — The watchman stays awake all night.
  \item \emph{Recall:} say the Sanskrit for speed, then the Sanskrit for creeps up to, then say \emph{jāgarti} again
\end{itemize}

\begin{tcolorbox}[breakable,colback=teal!4,colframe=teal!35!black,title={Before you move on}]
What does \sk{जागर्ति} mean? (\textquotedblleft{}Stays awake; keeps watch\textquotedblright{}.) Say it once more.
\end{tcolorbox}
\section[\texorpdfstring{\sk{कोलाहलः} (kolāhalaḥ)}{कोलाहलः (kolāhalaḥ)}]{\sk{कोलाहलः} (kolāhalaḥ) — an uproar}

\label{lesson:SA-C287-kolahalah}



\begin{quote}
Before the new one: say the Sanskrit for creeps up to, then the Sanskrit for stays awake.
\end{quote}

\subsection*{You'll want to know: \sk{कोलाहलः}}
\textbf{\sk{कोलाहलः}} — \emph{kolāhalaḥ} — \textquotedblleft{}an uproar\textquotedblright{}.

\begin{quote}
\sk{यदि} \sk{कोलाहलः} \sk{अभविष्यत्}, \sk{तर्हि} \sk{यामिकः} \sk{आगमिष्यत्।}
\end{quote}

\begin{quote}
\emph{yadi kolāhalaḥ abhaviṣyat, tarhi yāmikaḥ āgamiṣyat.}
\end{quote}

\begin{quote}
If there had been an uproar, then the watchman would have come.
\end{quote}

A masculine noun.

\begin{cousinweb}
\textbf{\sk{कोलाहलः}} (\emph{kolāhalaḥ}) comes from the sound itself: MW marks \emph{kolāhala} as onomatopoeic, a loud and confused noise of people or animals (MW).
\end{cousinweb}

\begin{grammarlens}[title={What you've built}]
One more word for telling what happened.
\end{grammarlens}

\subsection*{Your turn}
\begin{itemize}
\raggedright
  \item \emph{Say it:} \emph{kolāhalaḥ}
  \item \emph{Say it:} \emph{kolāhalaḥ}, once more
  \item \emph{Say it:} \emph{yadi kolāhalaḥ abhaviṣyat, tarhi yāmikaḥ āgamiṣyat.} — If there had been an uproar, then the watchman would have come.
  \item \emph{Recall:} say the Sanskrit for creeps up to, then the Sanskrit for stays awake, then say \emph{kolāhalaḥ} again
\end{itemize}

\begin{tcolorbox}[breakable,colback=teal!4,colframe=teal!35!black,title={Before you move on}]
What does \sk{कोलाहलः} mean? (\textquotedblleft{}An uproar\textquotedblright{}.) Say it once more.
\end{tcolorbox}
